% !TeX root = ../main.tex
% Add the above to each chapter to make compiling the PDF easier in some editors.

\chapter{Background and Related Work}\label{chapter:Related_Work}

This chapter explains how this thesis relates to the previous research in the fields it builds on: document extraction, biomedical information extraction, and LLM-based pipelines. The components it requires -- PDF and table extractors, biomedical entity normalizers, LLM-based structured extraction, natural language inference, and cost-aware model cascades -- are each relatively mature in isolation, but are rarely integrated into a single auditable, end-to-end system whose quality is evaluated under cost and accuracy constraints. Each section explains what the thesis borrows from prior work and where its design differs.

This chapter follows the structure of the system it describes. Sections~\ref{section:background-scientific-doc-processing}--\ref{section:background-table-and-figure-extraction} cover document extraction: how scientific documents are given a logical document structure, the PDF and layout-aware tools used by the pipeline, and the extraction of tables and figures. Sections~\ref{section:background-biomedical-info-extraction} and~\ref{section:background-biomedical-entity-normalization-and-resources} discuss biomedical information extraction and the entity normalization process that is used to make findings comparable across papers. Sections~\ref{section:background-llm-based-structured-extraction-and-summarization} and~\ref{section:background-grounding-entailment-and-nli} cover LLM-based structured knowledge extraction and natural language inference, which are used to extract claims, check whether they are supported by the evidence they cite, and classify relations between extracted rules. Sections~\ref{section:background-provenance-traceability-and-auditability} and~\ref{section:background-cost-aware-llm-inference-and-cascading} discuss the provenance-preservation requirement of the system and the cost-awareness requirement of the cascade. Section~\ref{section:background-evaluation-of-llm-extraction-systems} discusses how LLM-based extraction systems can be evaluated when the labels themselves are LLM-generated.
Finally, Section~\ref{section:background-positioning-of-thesis} brings all of these sections together and clarifies the thesis's contribution.

\section{Scientific document processing}\label{section:background-scientific-doc-processing}

The thesis operates on histopathology papers in English language, downloaded from PubMed Central Open Access subset~\parencite{nlm_pmc_open_access}. PDF is a common distribution format used by many articles in scientific literature, which is designed for visual representation, rather than to convey a logical structure; therefore, the section hierarchy, reading order, tables, figures, captions, and the relationship between all of these should be reconstructed~\parencite{}.

This reconstruction is especially important in a biomedical setting: clinically relevant statements and claims are not confined to running text; they may appear in tables, figures, captions or footnotes. Therefore, a pipeline that uses only body text risks discarding part of the evidence it is later asked to cite and ground~\parencite{}. For this reason, the pipeline implemented in this system treats the recovery of text, sections, figures, tables, captions, footnotes, and the bounding box coordinates of each of these artifacts as a first-class stage (Section~\ref{section:Document_Extraction}) rather than a post-processing stage as an afterthought; then the pipeline carries these coordinates forward to later stages, so that every later claim can be traced back to its source, which can be located on the specific article. 
\section{PDF and layout-aware extraction tools}\label{section:background-pdf-and-layout-aware-extraction-tools}

Several tools exist for converting scientific PDFs into structured representations, ranging from rule-based parsers to layout-aware neural models, and recent document-conversion systems that work end-to-end~\parencite{}. This thesis uses Docling as the main structural parser backbone of the document-extraction stage~\parencite{}. Docling is suitable for this role, because it outputs typed objects, such as text blocks, headings, figures, and tables, each with their bounding boxes. A section hierarchy and the reading order are reconstructed from this output. 

Docling serves as the starting point for the pipeline, whose output quality is improved by adding post-processing stages, such as a hybrid table detector, a two-pass filter for invisible text, a caption-matcher, and a crop-geometry correction mechanism, all of which are described in Section~\ref{section:Document_Extraction}. Then, the contribution of these additions is measured, and compared to an off-the-shelf Docling baseline in Chapter~\ref{chapter:Results}. Several other converters were examined, including Nougat, Marker, PDFFigures2, and PDFMiner; but these alternatives were not selected for production~\parencite{}.
\section{Table and figure extraction}\label{section:background-table-and-figure-extraction}

Figures, and especially tables, are the parts of a paper that are the most difficult for a layout parser to recover. However, the central quantitative evidence is presented in such formats in histopathology papers. Table extraction is commonly divided into two separate tasks: table detection, and table structure recognition. Table detection involves locating the tables and their boundaries, whereas table structure recognition involves the reconstruction of rows, columns, and cells in the table. Recent work on these tasks increasingly uses image-based transformer detectors~\parencite{}. Figure extraction and the association of captions to figures follow a similar line, with systems designed to detect figures, extract their captions, and associate them correctly~\parencite{}.

The thesis follows the logic in recent work, separating the steps of detection, structure recognition, and association. The document-extraction stage combines a layout-based detector (Docling) with an image-based table detector (Microsoft's Table Transformer) in a hybrid detector, and adds footnote-expansion, geometric caption-matching, and decorative artifact filtering to the system, as described in Section~\ref{section:Document_Extraction}. The results in Chapter~\ref{chapter:Results} show that off-the-shelf Docling already localizes tables and figures well; the improvements over the off-the-shelf Docling baseline mostly come from footnote capture for tables, and from caption attachment and decorative-icon filtering for figures.
\section{Biomedical information extraction}\label{section:background-biomedical-info-extraction}
\section{Biomedical entity normalization and resources}\label{section:background-biomedical-entity-normalization-and-resources}

Papers use different terminology for the same concepts, so a normalization process is necessary for the pipeline to compare findings across papers. This is the task of biomedical entity normalization: it maps surface forms that occur in papers' text to a controlled vocabulary. This thesis combines scispaCy biomedical models~\parencite{neumann2019scispacy} with the Unified Medical Language System (UMLS)~\parencite{bodenreider2004umls} to assign concept identifiers (CUIs) to subject and outcome entities, as described in Section~\ref{section:Knowledge_Extraction}.

The assigned CUIs provide the comparability gate with a stable way to determine whether two findings refer to the same entity. If both findings have a CUI, the gate compares those CUIs directly; otherwise, normalized-string matching is used as a fallback. The later grouping and canonicalization stages do not depend on CUI assignment; instead, they use the normalized subject and outcome strings produced by the same normalization step. The UMLS-derived entity buckets also contribute to the ILP-based calibration-cluster selection (§\ref{subsec:eval-ilp}), where CUI overlap is one of several weighted relatedness signals. Alternative biomedical linkers and an evaluation benchmark exist~\parencite{}, but the scispaCy/UMLS stack was found sufficient for the needs of this thesis. 
\section{LLM-based structured extraction and summarization}\label{section:background-llm-based-structured-extraction-and-summarization}

Large language models are increasingly used to read text and emit structured outputs under a predefined schema, which is the mechanism used in the MAP stage of this thesis. Two properties of LLM-based extraction are important for the pipeline design.

First, a single LLM call is not reliable enough by itself, and its output should be checked, as it may hallucinate, paraphrase evidence too freely, or produce different outputs across runs~\parencite{huang2025surveyhallucinationlarge}. For these reasons, agreement-based strategies such as self-consistency are used in the literature, where multiple output generations are collected and reconciled into a more reliable output~\parencite{wang2022selfconsistency}. This thesis generalizes this idea by extending repeated output generation from the same model to cross-provider voting, via the agreement-based cascading mechanism of Section~\ref{section:Agreement_Based_Cascading}. Other reliability techniques are not used here. Chain-of-thought prompting can make multi-step extraction more reliable by introducing intermediate reasoning steps~\parencite{wei2022chainofthoughtpromptingelicits}, but the pipeline does not treat these reasoning traces as evidence; instead, it only retains the verbatim source spans stored with each finding to support grounding and provenance (§\ref{section:Knowledge_Extraction}). A related line of work repairs the unreliable LLM outputs, either by using a different LLM to correct the existing output~\parencite{zhong2024harnessinglargelanguage}, or by editing the claim until a fact-verification model scores it as supported~\parencite{chen2023convergetruthfactual}. In contrast, this thesis uses agreement across providers to reconcile outputs and drops unsupported claims instead of rewriting them. 

Second, the thesis favors typed, atomic findings over summarization-style consolidation of free-form text: atomic claims are easily groundable and auditable, whereas free-text summaries merge multiple claims and make it harder to trace claims back to their evidence, as detailed in Section~\ref{section:Knowledge_Extraction}.
\section{Grounding, entailment, and NLI}\label{section:background-grounding-entailment-and-nli}

Checking whether a claim is supported by the evidence text it cites can be formulated as a natural-language-inference (NLI) problem: does the evidence paragraph entail the claim? In this thesis, this NLI problem is answered with a frozen biomedical NLI cross-encoder model, based on PubMedBERT~\parencite{gu2021biomedbert} and fine-tuned on the MNLI~\parencite{williams2018mnli} and MedNLI~\parencite{romanov2018mednli} entailment datasets, which build on the broader NLI tradition represented by SNLI~\parencite{}. This NLI model is applied following the sentence-pair scoring setup used in related representation-learning work~\parencite{reimers2019sbert}. 

The same NLI model is used in two different roles in the thesis pipeline: as a grounding filter, and as a relation classifier. The grounding filter scores each (claim, evidence text) pair and drops findings that are not supported enough by their evidence; the relation classifier scores the entailment and contradiction of each rule pair, and assigns a supporting, contradicting, or unrelated label to each of these pairs; the logic used for the label assignment is described in Section~\ref{section:Knowledge_Extraction}. The use of entailment scores as a faithfulness check connects the thesis to other work on detecting unsupported or hallucinated content by comparing generated claims against evidence~\parencite{}. 

The contribution of this thesis is the empirical results obtained from using an off-the-shelf biomedical NLI model with frozen thresholds, showing that such use of these models performs sufficiently well for both source faithfulness and relation classification. The obtained empirical results for both roles of the NLI model are reported in Chapter~\ref{chapter:Results}. 
\section{Provenance, traceability, and auditability}\label{section:background-provenance-traceability-and-auditability}
\section{Cost-aware LLM inference and cascading}\label{section:background-cost-aware-llm-inference-and-cascading}
\section{Evaluation of LLM extraction systems}\label{section:background-evaluation-of-llm-extraction-systems}

The reference labels are not human-annotated, but are also generated by an LLM. When human gold annotations are not available, recent work uses silver labels generated by strong models as a scalable alternative, but this introduces a clear risk: an evaluation based on silver labels can inherit the biases, preferences, or failure modes of the model that generates the labels~\parencite{}.

This thesis uses a human-labeled rubric for document extraction, and Opus-generated silver labels for the scoring and evaluation of knowledge extraction. The construction of these silver labels, the implications of measuring agreement with a silver labeler from a shared model family instead of clinical truth, and the offline-replay protocol and the cost-quality frontier methodology that keep experiments reproducible and the trade-off visible are deferred to Chapters~\ref{chapter:Methodology} and~\ref{chapter:Discussion}.
\section{Positioning of this thesis}\label{section:background-positioning-of-thesis}

The preceding sections describe components that each have their own literature, but are usually treated separately: table, figure, and text extraction from PDFs, biomedical entity normalization, LLM-based extraction, natural-language inference, and cost-aware model routing. The position of this thesis is that an auditable clinical knowledge base cannot be built by any of these components alone, and should be integrated into a single pipeline, in which provenance is preserved, every artifact, and more importantly, every emitted rule can be traced back to its source, and finally, quality is evaluated together with cost over the whole pipeline. 

The system built in this thesis, therefore, consists of three parts. First, it uses a provenance-preserving document extraction stage, built on Docling, and extended with a hybrid table detector and post-processing stages. Second, it implements a grounded, atomic, and normalized rule-extraction pipeline, preserving source-faithfulness, and using scispaCy and UMLS with a frozen biomedical NLI backbone~\parencite{}. Third, it adds an agreement-based cascading pipeline~\parencite{} that is calibrated on a histopathology corpus. 

The main contribution of the thesis is the combination of all these components: the thesis measures what it costs to keep every extracted rule traceable to its evidence, where existing components are already sufficient off-the-shelf, and where additional stages in the pipeline, such as grounding, agreement, and escalation, are worth their cost. The results are reported and discussed in Chapters~\ref{chapter:Results} and~\ref{chapter:Discussion}.